% defer/rcuusage.tex
% mainfile: ../perfbook.tex
% SPDX-License-Identifier: CC-BY-SA-3.0

\subsection{RCU的使用}
\label{sec:defer:RCU Usage}
\OriginallyPublished{sec:defer:RCU Usage}{RCU Usage}{Linux Weekly News}{PaulEMcKenney2008WhatIsRCUUsage}

本节从RCU可被应用的场景这一角度来回答``什么是RCU？''这个问题。
由于RCU最常被用来替换某些既有机制，
我们主要从它与这些机制的关系来审视它，
如\cref{tab:defer:RCU Usage}所列
以及\cref{fig:defer:Relationships Between RCU Use Cases}所展示的那样。
在本表所列各节之后，
\cref{sec:defer:RCU Usage Summary}给出了总结，
并在后续各节和其他文献~\cite{PaulEMcKenney2021RCUMysteriesFundamental,PaulEMcKenney2022RCUMysteriesUseCases}中进一步展开。

\begin{table}
\renewcommand*{\arraystretch}{1.2}
\centering
\small
\begin{tabular}{ll}
\toprule
RCU所替换的机制 & 页码 \\
\midrule
用于pre-BSD路由的RCU &
	\pageref{sec:defer:RCU for Pre-BSD Routing} \\
等待已有事务完成 &
	\pageref{sec:defer:Wait for Pre-Existing Things to Finish} \\
分阶段状态变更 &
	\pageref{sec:defer:Phased State Change} \\
仅添加链表（发布/订阅） &
	\pageref{sec:defer:Add-Only List} \\
类型安全内存 &
	\pageref{sec:defer:Type-Safe Memory} \\
存在性保证 &
	\pageref{sec:defer:Existence Guarantee} \\
轻量级垃圾收集器 &
	\pageref{sec:defer:Light-Weight Garbage Collector} \\
仅删除链表 &
	\pageref{sec:defer:Delete-Only List} \\
准读写锁 &
	\pageref{sec:defer:Quasi Reader-Writer Lock} \\
准引用计数 &
	\pageref{sec:defer:Quasi Reference Count} \\
准多版本并发控制（MVCC） &
	\pageref{sec:defer:Quasi Multi-Version Concurrency Control} \\
\bottomrule
\end{tabular}
\caption{RCU的使用}
\label{tab:defer:RCU Usage}
\end{table}

\subsubsection{用于Pre-BSD路由的RCU}
\label{sec:defer:RCU for Pre-BSD Routing}

与后面各节不同，本节聚焦于一个非常具体的用例，以便与其他机制进行比较。

\Cref{lst:defer:RCU Pre-BSD Routing Table Lookup,%
lst:defer:RCU Pre-BSD Routing Table Add/Delete}
展示了受RCU保护的Pre-BSD路由表代码
（\path{route_rcu.c}）。
前者展示了数据结构和\co{route_lookup()}，
后者展示了\co{route_add()}和\co{route_del()}。

\begin{listing}
\input{CodeSamples/defer/route_rcu@lookup.fcv}
\caption{RCU Pre-BSD路由表查找}
\label{lst:defer:RCU Pre-BSD Routing Table Lookup}
\end{listing}

\begin{listing}
\input{CodeSamples/defer/route_rcu@add_del.fcv}
\caption{RCU Pre-BSD路由表添加/删除}
\label{lst:defer:RCU Pre-BSD Routing Table Add/Delete}
\end{listing}

\begin{fcvref}[ln:defer:route_rcu:lookup]
在\cref{lst:defer:RCU Pre-BSD Routing Table Lookup}中，
\clnref{rh}添加了RCU回收所用的\co{->rh}字段，
\clnref{re_freed}添加了释放后使用检查字段\co{->re_freed}，
\clnref{lock,unlock1,unlock2}
添加了RCU读侧保护，
而\clnref{chk_freed,abort}则添加了释放后使用检查。
\end{fcvref}
\begin{fcvref}[ln:defer:route_rcu:add_del]
在\cref{lst:defer:RCU Pre-BSD Routing Table Add/Delete}中，
\clnref{add:lock,add:unlock,del:lock,%
del:unlock1,del:unlock2}添加了更新侧加锁，
\clnref{add:add_rcu,del:del_rcu}添加了RCU更新侧保护，
\clnref{del:call_rcu}使得\co{route_cb()}在
grace period过后被调用，
而\clnrefrange{cb:b}{cb:e}则定义了\co{route_cb()}。
对于一个可工作的并发实现来说，这已是最少量的附加代码。
\end{fcvref}

\begin{figure}
\centering
\resizebox{2.5in}{!}{\includegraphics{CodeSamples/defer/data/hps.2019.12.17a/perf-rcu}}
\caption{受RCU保护的Pre-BSD路由表}
\label{fig:defer:Pre-BSD Routing Table Protected by RCU}
\end{figure}

\Cref{fig:defer:Pre-BSD Routing Table Protected by RCU}
展示了在只读工作负载下的性能。
RCU的可扩展性相当好，提供了近乎理想的性能。
然而，这些数据是使用用户空间
RCU~\cite{MathieuDesnoyers2009URCU,PaulMcKenney2013LWNURCU}的\co{RCU_SIGNAL}
风格生成的，
其中\co{rcu_read_lock()}和\co{rcu_read_unlock()}
会生成少量代码。
对于\IXacr{qsbr}风格的RCU——\co{rcu_read_lock()}和\co{rcu_read_unlock()}
完全不生成任何代码——会发生什么呢？
（参见\cref{sec:defer:Introduction to RCU}，
特别是
\cref{fig:defer:QSBR: Waiting for Pre-Existing Readers}，
其中讨论了RCU QSBR。）

答案如
\cref{fig:defer:Pre-BSD Routing Table Protected by RCU QSBR}所示，
这表明RCU QSBR的性能和可扩展性实际上超过了
理想的无同步工作负载。

\QuickQuizSeries{%
\QuickQuizB{
	等等，什么？？？
	RCU QSBR怎么可能比理想情况还好？
	究竟是什么荒唐的"理想"定义才会不是
	所有可能结果中最好的？？？
}\QuickQuizAnswerB{
	这是一个极好的问题，答案是现代CPU和编译器极其复杂。
	但在深入探讨之前，非常值得注意的是，
	RCU QSBR的性能优势仅出现在
	每核一个硬件线程的情况下。
	一旦系统满载，RCU QSBR的性能就会
	回落到理想水平。

	\co{route_lookup()}搜索循环的RCU变体实际上
	比顺序版本多出一条x86指令，
	即序列\co{cmp}、\co{je}、\co{mov}、\co{cmp}、\co{lea}和\co{jne}中的\co{lea}。
	这条额外的指令是由于RCU变体的\co{route_entry}结构
	开头的\co{rcu_head}结构所致，
	因此，与顺序变体不同，RCU变体的
	\co{->re_next.next}指针有一个非零偏移量。
	回到20世纪80年代，这条额外的\co{lea}指令
	可能会可靠地导致RCU变体更慢，但我们
	现在已经是21\textsuperscript{世纪}了，80年代
	早已过去。

	但那些仔细阅读了
	\cref{sec:cpu:Pipelined CPUs}
	的读者早已知道这一切！

	这些反直觉的结果当然意味着，在现代微处理器上
	任何性能结果都必须持一定的怀疑态度。
	从理论上讲，获得优于理想的性能结果确实没有道理，
	但在现代微处理器上这确实可能发生。
	这样的结果可以被看作类似于著名的
	超线性加速（参见
	\cref{sec:SMPdesign:Beyond Partitioning}
	中的一个例子），即有趣但实际重要性有限。
	尽管如此，RCU的一大优势就是其读侧
	开销如此之低，以至于像这样的微小效应
	在实际的性能测量中清晰可见。

\begin{figure}
\centering
% Run with initial rcu_head structure in route_entry moved down.
\resizebox{2.5in}{!}{\includegraphics{CodeSamples/defer/data/hps.2019.12.17a/perf-rcu-qsbr}}
\caption{受RCU QSBR保护、\tco{rcu_head}非首字段的Pre-BSD路由表}
\label{fig:defer:Pre-BSD Routing Table Protected by RCU QSBR With Non-Initial rcu-head}
\end{figure}

	这就引出了一个问题：如果将
	\co{rcu_head}结构移动使得RCU的
	\co{->re_next.next}指针也具有零偏移量，
	与顺序变体完全相同，会发生什么。
	答案如
	\cref{fig:defer:Pre-BSD Routing Table Protected by RCU QSBR With Non-Initial rcu-head}
	所示，这会导致RCU QSBR的性能下降到
	仍然非常接近理想的水平，但不再是超理想的。
}\QuickQuizEndB
%
\QuickQuizE{
	鉴于RCU QSBR的读侧性能如此出色，为何还要费心使用
	用户空间RCU的其他风格呢？
}\QuickQuizAnswerE{
	因为RCU QSBR对整个应用程序施加了一些
	可能无法容忍的约束，
	例如，要求应用程序中的每一个线程
	都定期通过一个quiescent state。
	除此之外，这意味着RCU QSBR对
	库的编写者没有帮助，他们可能更适合使用其他
	风格的用户空间RCU~\cite{PaulMcKenney2013LWNURCU}。
}\QuickQuizEndE
}

\begin{figure}
\centering
% Run with initial rcu_head structure at beginning of route_entry structure.
% This need to run twice is the reason for the oddball directory.
% @@@ Make this generate both files with a single run?
\resizebox{2.5in}{!}{\includegraphics{CodeSamples/defer/data/hps.2019.12.02a/perf-rcu-qsbr}}
\caption{受RCU QSBR保护的Pre-BSD路由表}
\label{fig:defer:Pre-BSD Routing Table Protected by RCU QSBR}
\end{figure}

\begin{figure*}
\centering
\IfTwoColumn{
  \resizebox{5.5in}{!}{\includegraphics{defer/RCUusecases}}
}{
  \resizebox{.96\textwidth}{!}{\includegraphics{defer/RCUusecases}}
  % eb builds require .96
}
\caption{RCU各用例之间的关系}
\label{fig:defer:Relationships Between RCU Use Cases}
\end{figure*}

虽然Pre-BSD路由是一个极好的RCU用例，但还是值得
看看
\cref{fig:defer:Relationships Between RCU Use Cases}
所示的更广泛的用例之间的关系。
以下各节将承担这一任务。

在阅读这些节的同时，请思考哪个用例
最能描述Pre-BSD路由。

\subsubsection{等待已有事务完成}
\label{sec:defer:Wait for Pre-Existing Things to Finish}

如\cref{sec:defer:RCU Fundamentals}所述，
RCU的一个重要组成部分
是等待RCU读者完成的方式。
RCU的一大优势是它允许你等待数千种不同的事务完成，
而无需显式地跟踪每一个，也不会招致
使用显式跟踪的方案所固有的
性能下降、可扩展性限制、复杂的死锁
场景和内存泄漏风险。

在本节中，我们将展示\co{synchronize_sched()}的
读侧对应物（包括任何禁用抢占的操作，
以及硬件操作和
禁用中断的原语）如何允许你与
不可屏蔽中断（NMI）处理程序交互，这用锁机制来做是相当困难的。
这种方法被称为``纯RCU''~\cite{PaulEdwardMcKenneyPhD}，
在Linux内核中的少数地方有使用。

这种``纯RCU''设计的基本形式如下：

\begin{enumerate}
\item	进行一项变更，例如改变操作系统对NMI的响应方式。
\item	等待所有已有的read-side critical section
	完全结束（例如，使用
	\co{synchronize_sched()}原语）。\footnote{
		在Linux内核v5.1及更高版本中，\co{synchronize_sched()}已
		被纳入\co{synchronize_rcu()}。}
	这里的关键观察是，后续的RCU read-side critical section
	保证能看到所做的任何变更。
\item	清理，例如返回表示变更已成功完成的状态。
\end{enumerate}

本节的其余部分展示了改编自Linux内核的示例代码。
在此示例中，现已废弃的oprofile设施中的
\co{nmi_stop()}函数使用\co{synchronize_sched()}来确保
所有进行中的NMI通知在释放相关
资源之前已经完成。
此代码的简化版本如
\cref{lst:defer:Using RCU to Wait for NMIs to Finish}所示。

\begin{listing}
\begin{fcvlabel}[ln:defer:Using RCU to Wait for NMIs to Finish]
\begin{VerbatimL}[commandchars=\\\@\$]
struct profile_buffer {				\lnlbl@struct:b$
	long size;
	atomic_t entry[0];
};						\lnlbl@struct:e$
static struct profile_buffer *buf = NULL;	\lnlbl@struct:buf$

void nmi_profile(unsigned long pcvalue)		\lnlbl@nmi_profile:b$
{
	struct profile_buffer *p = rcu_dereference(buf);\lnlbl@nmi_profile:rcu_deref$

	if (p == NULL)				\lnlbl@nmi_profile:if_NULL$
		return;				\lnlbl@nmi_profile:ret:a$
	if (pcvalue >= p->size)			\lnlbl@nmi_profile:if_oor$
		return;				\lnlbl@nmi_profile:ret:b$
	atomic_inc(&p->entry[pcvalue]);		\lnlbl@nmi_profile:inc$
}						\lnlbl@nmi_profile:e$

void nmi_stop(void)				\lnlbl@nmi_stop:b$
{
	struct profile_buffer *p = buf;		\lnlbl@nmi_stop:fetch$

	if (p == NULL)				\lnlbl@nmi_stop:if_NULL$
		return;				\lnlbl@nmi_stop:ret$
	rcu_assign_pointer(buf, NULL);		\lnlbl@nmi_stop:NULL$
	synchronize_sched();			\lnlbl@nmi_stop:sync_sched$
	kfree(p);				\lnlbl@nmi_stop:kfree$
}						\lnlbl@nmi_stop:e$
\end{VerbatimL}
\end{fcvlabel}
\caption{使用RCU等待NMI完成}
\label{lst:defer:Using RCU to Wait for NMIs to Finish}
\end{listing}

\begin{fcvref}[ln:defer:Using RCU to Wait for NMIs to Finish:struct]
\Clnrefrange{b}{e}定义了一个\co{profile_buffer}结构，包含
一个大小字段和一个不定长的条目数组。
\Clnref{buf}定义了一个指向profile缓冲区的指针，
该指针大概在其他地方被初始化为指向一块动态分配的
内存区域。
\end{fcvref}

\begin{fcvref}[ln:defer:Using RCU to Wait for NMIs to Finish:nmi_profile]
\Clnrefrange{b}{e}定义了\co{nmi_profile()}函数，
它在NMI处理程序内被调用。
因此，它不能被抢占，也不能被普通的
中断处理程序中断，然而，它仍然可能由于缓存未命中、
ECC错误以及同一核心内其他硬件线程的周期窃取而延迟。
\Clnref{rcu_deref}使用\co{rcu_dereference()}原语获取
指向profile缓冲区的本地指针，以确保在
DEC Alpha上的内存排序，
\clnref{if_NULL,ret:a}在没有当前分配的
profile缓冲区时从此函数退出，而\clnref{if_oor,ret:b}
在\co{pcvalue}参数超出范围时
从此函数退出。
否则，\clnref{inc}递增以\co{pcvalue}参数为索引的
profile缓冲区条目。
注意，将大小与缓冲区一起存储可以保证
范围检查与缓冲区匹配，即使一个大缓冲区突然
被一个较小的缓冲区替换。
\end{fcvref}

\begin{fcvref}[ln:defer:Using RCU to Wait for NMIs to Finish:nmi_stop]
\Clnrefrange{b}{e}定义了\co{nmi_stop()}函数，
调用者负责互斥（例如，持有正确的锁）。
\Clnref{fetch}获取指向profile缓冲区的指针，
\clnref{if_NULL,ret}在没有缓冲区时退出函数。
否则，\clnref{NULL}将profile缓冲区指针置\co{NULL}
（使用\co{rcu_assign_pointer()}原语在弱序
机器上维持内存排序），
\clnref{sync_sched}等待一个RCU Sched grace period过去，
特别是等待所有不可抢占的代码区域
（包括NMI处理程序）完成。
一旦执行继续到\clnref{kfree}，我们就可以保证
任何获取了旧缓冲区指针的\co{nmi_profile()}
实例已经返回。
因此，可以安全地释放缓冲区，本例中使用的是
\co{kfree()}原语。
\end{fcvref}

\QuickQuiz{
	假设\co{nmi_profile()}函数是可抢占的。
	要使此示例正确工作，需要做哪些改变？
}\QuickQuizAnswer{
	一种方法是在\co{nmi_profile()}中使用
	\co{rcu_read_lock()}和\co{rcu_read_unlock()}，
	并将\co{synchronize_sched()}替换为\co{synchronize_rcu()}，
	也许如
	\cref{lst:defer:Using RCU to Wait for Mythical Preemptible NMIs to Finish}所示。
%
\begin{listing}
\begin{VerbatimL}
struct profile_buffer {
	long size;
	atomic_t entry[0];
};
static struct profile_buffer *buf = NULL;

void nmi_profile(unsigned long pcvalue)
{
	struct profile_buffer *p;

	rcu_read_lock();
	p = rcu_dereference(buf);
	if (p == NULL) {
		rcu_read_unlock();
		return;
	}
	if (pcvalue >= p->size) {
		rcu_read_unlock();
		return;
	}
	atomic_inc(&p->entry[pcvalue]);
	rcu_read_unlock();
}

void nmi_stop(void)
{
	struct profile_buffer *p = buf;

	if (p == NULL)
		return;
	rcu_assign_pointer(buf, NULL);
	synchronize_rcu();
	kfree(p);
}
\end{VerbatimL}
\caption{使用RCU等待虚构的可抢占NMI完成}
\label{lst:defer:Using RCU to Wait for Mythical Preemptible NMIs to Finish}
\end{listing}

	但NMI处理程序为什么会是可抢占的呢？？？
}\QuickQuizEnd

简而言之，RCU使得动态切换profile缓冲区变得容易
（你不妨\emph{试试}用原子操作高效地做到这一点，
或者用锁来做！）。
这是RCU纯粹形式的一种罕见使用。
RCU通常在更高的抽象层次上使用，
如后续各节所示。

\subsubsection{分阶段状态变更}
\label{sec:defer:Phased State Change}

\Cref{fig:defer:Phased State Change for Maintenance Operation}
展示了一个分阶段状态变更的时间线示例，
用于高效地处理维护操作。
如果没有正在进行的维护操作，常规操作
必须快速执行，例如不需要获取读写锁。
然而，如果有维护操作正在进行，常规操作
必须谨慎进行，考虑到因与维护操作
并发运行而增加的复杂性。
这意味着在维护操作期间，常规操作将招致更高的开销，
这也是维护操作通常被安排在
低负载时段执行的原因之一。

\begin{figure}
\centering
\resizebox{\twocolumnwidth}{!}{\includegraphics{defer/RCUphasedstatechange}}
\caption{维护操作的分阶段状态变更}
\label{fig:defer:Phased State Change for Maintenance Operation}
\end{figure}

在图中，这些看似矛盾的需求通过
在维护操作之前设置准备阶段、在之后设置清理阶段来解决，
在这些阶段中常规操作可以快速执行或谨慎执行。

\begin{listing}
\begin{fcvlabel}[ln:defer:Phased State Change for Maintenance Operations]
\begin{VerbatimL}[commandchars=\\\@\$]
bool be_careful;

void cco(void)
{
	rcu_read_lock();			\lnlbl@rrl$
	if (READ_ONCE(be_careful))		\lnlbl@if$
		cco_carefully();		\lnlbl@careful$
	else
		cco_quickly();			\lnlbl@quick$
	rcu_read_unlock();			\lnlbl@rul$
}

void maint(void)
{
	WRITE_ONCE(be_careful, true);		\lnlbl@tocareful$
	synchronize_rcu();			\lnlbl@sr1$
	do_maint();				\lnlbl@maint$
	synchronize_rcu();			\lnlbl@sr2$
	WRITE_ONCE(be_careful, false);		\lnlbl@toquick$
}
\end{VerbatimL}
\end{fcvlabel}
\caption{维护操作的分阶段状态变更}
\label{lst:defer:Phased State Change for Maintenance Operations}
\end{listing}

\begin{fcvref}[ln:defer:Phased State Change for Maintenance Operations]
此分阶段状态变更的伪代码示例如
\cref{lst:defer:Phased State Change for Maintenance Operations}所示。
常规操作由\co{cco()}在从\clnref{rrl}到\clnref{rul}的
RCU read-side critical section内执行。
其中\clnref{if}检查全局的\co{be_careful}标志，
相应地调用\co{cco_carefully()}或\co{cco_quickly()}。
\end{fcvref}

\begin{fcvref}[ln:defer:Phased State Change for Maintenance Operations]
这允许\co{maint()}函数在\clnref{tocareful}设置\co{be_careful}标志，
并在\clnref{sr1}等待一个RCU grace period。
当控制到达\clnref{maint}时，所有看到\co{be_careful}
为\co{false}值（因此可能调用
\co{cco_quickly()}函数）的\co{cco()}函数将已经完成其操作，
从而所有当前正在执行的\co{cco()}函数都将调用
\co{cco_carefully()}。
这意味着此时调用\co{do_maint()}函数是安全的。
\Clnref{sr2}接着等待所有可能与\co{do_maint()}并发
运行的\co{cco()}函数完成，最后
\clnref{toquick}将\co{be_careful}标志重新设为\co{false}。
\end{fcvref}

\QuickQuizSeries{%
\QuickQuizB{
	\cref{lst:defer:Phased State Change for Maintenance Operations}中
	\co{maint()}函数第二次调用\co{synchronize_rcu()}
	的意义何在？
	在清理阶段让\co{cco()}调用\co{cco_carefully()}或
	\co{cco_quickly()}不是都可以吗？
}\QuickQuizAnswerB{
	问题在于\co{cco()}函数从\co{be_careful}的加载
	与\co{cco_quickly()}函数执行的任何内存加载之间
	没有排序保证。
	由于没有排序保证，如果没有第二次调用
	\co{synchronize_rcu()}，内存排序可能导致
	\co{cco_quickly()}中的加载与\co{do_maint()}中的存储重叠。

	另一种替代方案是通过将\co{READ_ONCE()}
	改为\co{smp_load_acquire()}、将
	\co{WRITE_ONCE()}改为\co{smp_store_release()}来
	弥补去掉第二次\co{synchronize_rcu()}调用的影响，
	从而恢复所需的排序。
}\QuickQuizEndB

\QuickQuizE{
	你怎么能确定
	\cref{lst:defer:Phased State Change for Maintenance Operations}中
	\co{maint()}所示的代码确实有效？
}\QuickQuizAnswerE{
	按照一种流行的学派观点，你无法确定。

	但在这个例子中，那些愿意跳到
	\cref{chp:Formal Verification}
	和
	\cref{chp:Advanced Synchronization: Memory Ordering}
	的读者可能会发现几个LKMM litmus测试很有趣
	（\path{C-RCU-phased-state-change-1.litmus}和
	\path{C-RCU-phased-state-change-2.litmus}）。
	可以说这些测试证明了这段代码
	及其一个变体确实能正常工作。
}\QuickQuizEndE
}% End of \QuickQuizSeries

分阶段状态变更允许频繁的操作使用轻量级检查，
无需昂贵的锁获取或原子读-修改-写操作，
并在Linux内核中以\co{rcu_sync}~\cite{OlegNesterov2013rcusync}的
形式被使用，以实现具有轻量级读者的
读写信号量的变体。
分阶段状态变更只是在等待完成用例
（\cref{sec:defer:Wait for Pre-Existing Things to Finish}）之上
增加了一个受检查的状态变量，
因此也处于较低的抽象层次。

\subsubsection{仅添加链表}
\label{sec:defer:Add-Only List}

仅添加数据结构，以仅添加链表为代表，可用于
数量惊人的常见用例，其中最常见的可能是
变更日志。
仅添加数据结构是对RCU底层
发布/订阅机制的纯粹使用。

Pre-BSD路由表的仅添加变体可以从
\cref{lst:defer:RCU Pre-BSD Routing Table Lookup,lst:defer:RCU Pre-BSD Routing Table Add/Delete}
派生而来。
由于没有删除操作，\co{route_del()}和\co{route_cb()}
函数可以省略，连同\co{route_entry}结构的
\co{->rh}和\co{->re_freed}字段、
\co{route_lookup()}函数中的
\co{rcu_read_lock()}和\co{rcu_read_unlock()}调用，
以及所有剩余函数中\co{->re_freed}字段的所有使用。

当然，如果有许多\co{route_add()}函数的并发调用，
\co{routelock}将面临严重的争用，
如果使用无锁技术，\co{routelist}也将面临严重的内存争用。
避免这种争用的通常方法是使用一种并发友好的
数据结构，如哈希表（参见\cref{chp:Data Structures}）。
另一种方法是将每CPU数据结构定期合并
到一个全局数据结构中。

另一方面，如果永远不进行删除操作，
长时间的多个\co{route_add()}并发调用最终将
耗尽所有可用内存。
因此，大多数受RCU保护的数据结构也实现了删除操作。

\subsubsection{类型安全内存}
\label{sec:defer:Type-Safe Memory}

许多无锁算法不要求给定的数据元素
在引用它的整个RCU read-side critical section期间
保持相同的身份——但前提是该数据元素保持
相同的类型。
换言之，这些无锁算法可以容忍给定数据元素
在被引用期间被释放并重新分配为相同类型的结构，
但必须禁止类型更改。
这种保证在学术文献中被称为``\IX{type-safe memory}''~\cite{Cheriton96a}，
它比\cref{sec:defer:Existence Guarantee}中讨论的
\IXpl{existence guarantee}更弱，
因此使用起来相当困难。
Linux内核中的类型安全内存算法利用slab缓存，
通过\co{SLAB_TYPESAFE_BY_RCU}特别标记这些缓存，
以便在将释放的slab返回到系统内存时使用RCU。
这种RCU的使用保证了这种slab中任何正在使用的元素
将在所有已有的RCU read-side critical section持续期间
保留在该slab中，从而保持其类型。

\QuickQuiz{
	但如果多个线程中有任意长的RCU read-side critical section
	序列，使得在任何时间点系统中总有至少一个线程
	在执行RCU read-side critical section呢？
	这不会阻止\co{SLAB_TYPESAFE_BY_RCU} slab中的
	任何数据被归还给系统，可能导致
	OOM事件吗？
}\QuickQuizAnswer{
	确实可能存在任意长的时间段，
	在此期间总有至少一个线程处于RCU read-side critical section中。
	然而，\cref{sec:defer:Type-Safe Memory}中
	描述的关键词是``正在使用''和``已有的''。
	请记住，给定的RCU read-side critical section
	在概念上只被允许获取对在该临界区期间
	对读者可见的数据元素的引用。
	此外，请记住slab不能被归还给系统，
	直到其所有数据元素都被释放——事实上，
	RCU grace period不能在它们全部被释放之前开始。

	因此，slab缓存只需等待那些在最后一个元素
	释放之前开始的RCU read-side critical section。
	这反过来意味着，在最后一个元素释放之后
	开始的任何RCU grace period都可以——在该grace period
	结束后，slab可以被归还给系统。
}\QuickQuizEnd

重要的是要注意，\co{SLAB_TYPESAFE_BY_RCU}
\emph{绝不会}阻止\co{kmem_cache_alloc()}立即
重新分配刚通过\co{kmem_cache_free()}释放的内存！
事实上，\co{rcu_dereference()}刚返回的受
\co{SLAB_TYPESAFE_BY_RCU}保护的数据结构
可能被释放和重新分配任意多次，
即使在\co{rcu_read_lock()}的保护下也是如此。
相反，\co{SLAB_TYPESAFE_BY_RCU}通过阻止
\co{kmem_cache_free()}将完全释放的数据结构slab
归还给系统来运作，直到一个RCU grace period过去之后。
简而言之，虽然给定的RCU read-side critical section
可能看到给定的\co{SLAB_TYPESAFE_BY_RCU}数据元素
被任意频繁地释放和重新分配，
但在该临界区完成之前，该元素的类型保证不会改变。

因此，这些算法通常使用一个验证步骤来检查
确保新引用的数据结构确实是
所请求的那个~\cite[Section~2.5]{LaninShasha1986TSM}。
这些验证检查要求数据结构的部分内容
在释放-重新分配过程中保持不变。
这样的验证检查通常很难做对，并且可能
隐藏微妙而难以发现的错误。

因此，虽然基于类型安全的无锁算法
在极少数困难情况下可以极为有用，
但你应当尽可能使用存在性保证。
毕竟，更简单几乎总是更好！
另一方面，基于类型安全的无锁算法可以
提供更好的缓存局部性，从而提高性能。
这种改善的缓存局部性源于这样一个事实：
此类算法可以立即重新分配刚释放的内存块。
相比之下，基于存在性保证的算法必须等待
所有已有的读者完成后才能重新分配内存，
届时该内存可能已被从CPU缓存中淘汰。

如\cref{fig:defer:Relationships Between RCU Use Cases}所示，
RCU的类型安全内存用例结合了等待完成
和发布/订阅两个组件，但在Linux内核中还包括
slab分配器通过\co{SLAB_TYPESAFE_BY_RCU}标志
指定的延迟回收。

\subsubsection{存在性保证}
\label{sec:defer:Existence Guarantee}

Gamsa等人~\cite{Gamsa99}
讨论了\IXpl{existence guarantee}，并描述了一种
类似RCU的机制如何能够提供这些存在性保证
（参见PDF第7页的第5节），
\cref{sec:locking:Lock-Based Existence Guarantees}
讨论了如何通过加锁来保证存在性，
以及这样做所带来的缺点。
其效果是，如果在RCU read-side critical section中
访问了任何受RCU保护的数据元素，
则该数据元素在整个RCU read-side critical section
期间都保证继续存在。

\begin{listing}
\begin{fcvlabel}[ln:defer:Existence Guarantees Enable Per-Element Locking]
\begin{VerbatimL}[commandchars=\\\@\$]
int delete(int key)
{
	struct element *p;
	int b;

	b = hashfunction(key);			\lnlbl@hash$
	rcu_read_lock();			\lnlbl@rdlock$
	p = rcu_dereference(hashtable[b]);
	if (p == NULL || p->key != key) {	\lnlbl@chkkey$
		rcu_read_unlock();		\lnlbl@rdunlock1$
		return 0;			\lnlbl@ret_0:a$
	}
	spin_lock(&p->lock);			\lnlbl@acq$
	if (hashtable[b] == p && p->key == key) {\lnlbl@chkkey2$
		rcu_read_unlock();		\lnlbl@rdunlock2$
		rcu_assign_pointer(hashtable[b], NULL);\lnlbl@remove$
		spin_unlock(&p->lock);		\lnlbl@rel1$
		synchronize_rcu();		\lnlbl@sync_rcu$
		kfree(p);			\lnlbl@kfree$
		return 1;			\lnlbl@ret_1$
	}
	spin_unlock(&p->lock);			\lnlbl@rel2$
	rcu_read_unlock();			\lnlbl@rdunlock3$
	return 0;				\lnlbl@ret_0:b$
}
\end{VerbatimL}
\end{fcvlabel}
\caption{存在性保证实现每元素加锁}
\label{lst:defer:Existence Guarantees Enable Per-Element Locking}
\end{listing}

\begin{fcvref}[ln:defer:Existence Guarantees Enable Per-Element Locking]
\Cref{lst:defer:Existence Guarantees Enable Per-Element Locking}
展示了基于RCU的存在性保证如何通过一个
从哈希表中删除元素的函数来实现每元素加锁。
\Clnref{hash}计算哈希函数，\clnref{rdlock}进入一个
RCU read-side critical section。
如果\clnref{chkkey}发现哈希表的对应桶
为空或其中的元素不是我们要删除的，
则\clnref{rdunlock1}退出RCU read-side critical section，
\clnref{ret_0:a}返回失败。
\end{fcvref}

\QuickQuiz{
	如果我们需要删除的元素不是
	\cref{lst:defer:Existence Guarantees Enable Per-Element Locking}
	\clnrefr{ln:defer:Existence Guarantees Enable Per-Element Locking:chkkey}行上
	链表的第一个元素怎么办？
}\QuickQuizAnswer{
	与（有缺陷的）
	\cref{lst:locking:Per-Element Locking Without Existence Guarantees (Buggy!)}
	一样，这是一个非常简单的无链接哈希表，因此给定桶中
	唯一的元素就是第一个元素。
	读者再次被邀请将此示例改编为具有
	完整链接的哈希表。
	不那么精力充沛的读者可以参阅
	\cref{chp:Data Structures}。
}\QuickQuizEnd

\begin{fcvref}[ln:defer:Existence Guarantees Enable Per-Element Locking]
否则，\clnref{acq}获取更新侧自旋锁，
\clnref{chkkey2}然后检查该元素是否仍然是我们想要的。
如果是，\clnref{rdunlock2}离开RCU read-side critical section，
\clnref{remove}将其从表中移除，\clnref{rel1}释放锁，
\clnref{sync_rcu}等待所有已有的RCU read-side critical section完成，
\clnref{kfree}释放新移除的元素，
\clnref{ret_1}返回成功。
如果该元素已不再是我们想要的，\clnref{rel2}释放锁，
\clnref{rdunlock3}离开RCU read-side critical section，
\clnref{ret_0:b}返回删除指定键失败。
\end{fcvref}

\QuickQuizSeries{%
\QuickQuizB{
	\begin{fcvref}[ln:defer:Existence Guarantees Enable Per-Element Locking]
	为什么在
	\cref{lst:defer:Existence Guarantees Enable Per-Element Locking}的
	\clnref{rdunlock2}行退出RCU read-side critical section，
	然后在\clnref{rel1}才释放锁，是可以的？
	\end{fcvref}
}\QuickQuizAnswerB{
	\begin{fcvref}[ln:defer:Existence Guarantees Enable Per-Element Locking]
	首先请注意，\clnref{chkkey2}上的第二次检查是必要的，
	因为在我们等待获取锁的期间，
	其他CPU可能已经移除了此元素。
	然而，我们在获取锁时处于RCU read-side critical section中，
	这保证了此元素不可能被重新分配
	并重新插入到此哈希表中。
	此外，一旦我们获取了锁，锁本身就保证了
	元素的存在，因此我们不再需要处于
	RCU read-side critical section中。
	% The question as to whether it is necessary to re-check the
	% element's key is left as an exercise to the reader.
	% A re-check is necessary if the key can mutate or if it is
	% necessary to reject deleted entries (in cases where deletion
	% is recorded by mutating the key.
	\end{fcvref}
}\QuickQuizEndB
%
\QuickQuizM{
	\begin{fcvref}[ln:defer:Existence Guarantees Enable Per-Element Locking]
	为什么不在
	\cref{lst:defer:Existence Guarantees Enable Per-Element Locking}的
	\clnref{rdunlock3}行退出RCU read-side critical section之前，
	先在\clnref{rel2}行释放锁呢？
	\end{fcvref}
}\QuickQuizAnswerM{
	假设我们颠倒这两行的顺序。
	那么这段代码就容易受到以下事件序列的影响：
	\begin{enumerate}
	\begin{fcvref}[ln:defer:Existence Guarantees Enable Per-Element Locking]
	\item	CPU~0调用\co{delete()}，找到了要删除的元素，
		执行到\clnref{rdunlock2}。
		它尚未实际删除该元素，但即将这样做。
	\item	CPU~1并发调用\co{delete()}，试图
		删除同一个元素。
		然而，CPU~0仍持有锁，因此CPU~1
		在\clnref{acq}处等待。
	\item	CPU~0执行\clnref{remove,rel1}，
		并在\clnref{sync_rcu}处阻塞，等待CPU~1
		退出其RCU read-side critical section。
	\item	CPU~1现在获取了锁，但\clnref{chkkey2}上的测试
		失败了，因为CPU~0已经移除了该元素。
		CPU~1现在执行\clnref{rel2}
		（为了本快速测验的目的，我们已将其与\clnref{rdunlock3}
		交换）
		并退出其RCU read-side critical section。
	\item	CPU~0现在可以从\co{synchronize_rcu()}返回，
		因此执行\clnref{kfree}，将元素发送到
		空闲列表。
	\item	CPU~1现在试图为一个已被释放的元素释放锁，
		更糟糕的是，该元素可能已被
		重新分配为其他类型的数据结构。
		这是一个致命的内存损坏错误。
	\end{fcvref}
	\end{enumerate}
}\QuickQuizEndM
%
\QuickQuizE{
	\cref{lst:defer:Existence Guarantees Enable Per-Element Locking}
	中所示的基于RCU的算法看起来与
	\cref{lst:locking:Per-Element Locking With Lock-Based Existence Guarantees}
	中的非常相似，那为什么基于RCU的方法应该更好呢？
}\QuickQuizAnswerE{
	\Cref{lst:defer:Existence Guarantees Enable Per-Element Locking}
	将
	\cref{lst:locking:Per-Element Locking With Lock-Based Existence Guarantees}
	中所示的每元素\co{spin_lock()}和\co{spin_unlock()}
	替换为开销低得多的\co{rcu_read_lock()}和\co{rcu_read_unlock()}，
	从而极大地提高了性能和可扩展性。
	更多详情请参阅
	\cref{sec:datastruct:RCU-Protected Hash Table Performance}。
}\QuickQuizEndE
}

细心的读者会认识到这只是原始的等待完成主题
（\cref{sec:defer:Wait for Pre-Existing Things to Finish}）
的一个轻微变体，添加了发布/订阅、链式结构、
堆分配器（通常）和延迟回收，如
\cref{fig:defer:Relationships Between RCU Use Cases}所示。
他们可能还会注意到，与
\cref{sec:locking:Lock-Based Existence Guarantees}
中讨论的基于锁的存在性保证相比，
这具有死锁免疫的优势。

\subsubsection{轻量级垃圾收集器}
\label{sec:defer:Light-Weight Garbage Collector}

初次了解RCU的人常发出的一个惊叹是``RCU有点像
垃圾收集器！''~\cite{MattKline2023LockBasedLocklessFresh}。
这一惊叹有很大的道理，特别是当使用RCU来
实现依赖垃圾收集的非阻塞算法时，
但有时也可能会产生误导。

也许思考RCU与自动垃圾收集器（GC）之间关系的最佳方式是：
RCU类似于GC，因为收集的\emph{时机}是自动
确定的，但RCU与GC的不同之处在于：
\begin{enumerate*}[(1)]
\item 程序员必须手动指明给定的数据结构
何时有资格被收集，且
\item 程序员必须手动标记可能持有引用的
RCU read-side critical section。
\end{enumerate*}

尽管存在这些差异，两者的相似性确实很深。
事实上，我所知的第一个类RCU机制使用了
基于引用计数的垃圾收集器来处理
grace period~\cite{Kung80}，RCU与垃圾收集之间的联系
也在较近期被注意到~\cite{HarshalSheth2016goRCU}。

轻量级垃圾收集器用例与存在性保证用例非常相似，
只是在混合中添加了所需的非阻塞算法。
此轻量级垃圾收集器用例也可以与下一节
描述的存在性保证结合使用。

\subsubsection{仅删除链表}
\label{sec:defer:Delete-Only List}

仅删除链表是\cref{sec:defer:Add-Only List}中
介绍的仅添加链表的较不常用的对应物，
可以被看作是没有发布/订阅组件的存在性保证用例，
如\cref{fig:defer:Relationships Between RCU Use Cases}所示。
仅删除链表可用于链表的所有可能成员
在初始化时已知、且成员可以被移除的情况。
例如，链表的元素可能代表系统中可能发生故障
但在重启前无法修复或替换的硬件组件。

Pre-BSD路由表的仅删除变体可以从
\cref{lst:defer:RCU Pre-BSD Routing Table Lookup,lst:defer:RCU Pre-BSD Routing Table Add/Delete}
派生而来。
由于没有添加操作，可以省略\co{route_add()}函数，
或者将其使用限制在初始化阶段。
理论上，\co{route_lookup()}函数可以使用非RCU的迭代器，
不过在Linux内核中这会引起调试代码的警告。
此外，RCU迭代器的增量成本通常可以忽略不计。

因此，仅删除的情况通常使用为
添加和删除而设计的算法和数据结构。

\subsubsection{准读写锁}
\label{sec:defer:Quasi Reader-Writer Lock}

也许Linux内核中最常见的RCU使用方式是
在读密集场景中替换读写锁。
然而，RCU的这种使用方式在一开始对我来说并不明显。
事实上，我在20世纪90年代早期
实现了一个轻量级读写锁~\cite{WilsonCHsieh92a}\footnote{
	类似于2.4 Linux内核中的\co{brlock}和
	更新版本Linux内核中的\co{lglock}。}，
然后才实现了通用的RCU实现。
我为轻量级读写锁设想的每一个使用场景
都改用了RCU来实现。
事实上，轻量级读写锁过了三年多才
迎来它的第一次使用。
真是让我大出洋相！

RCU与读写锁之间的关键相似之处在于
两者都有可以并发执行的读侧临界区。
事实上，在某些情况下，可以机械地将RCU API
成员替换为对应的读写锁API成员。
但首先，为什么要这样做？

RCU的优势包括性能、
死锁免疫和实时延迟。
当然，RCU也有局限性，包括
读者和更新者并发运行的事实、低优先级的RCU读者
可以阻塞等待grace period过去的高优先级线程、
以及grace period延迟可达多个毫秒。
以下段落将讨论这些优势和局限性。

\paragraph{性能}

\begin{figure}
\centering
\resizebox{2.5in}{!}{\includegraphics{CodeSamples/defer/data/rcuscale.hps.2020.05.28a/rwlockRCUperf}}
\caption{RCU相比读写锁的性能优势}
\label{fig:defer:Performance Advantage of RCU Over Reader-Writer Locking}
\end{figure}

\cref{fig:defer:Performance Advantage of RCU Over Reader-Writer Locking}
展示了Linux内核RCU相对于读写锁的读侧性能优势，
数据在一台448 CPU、2.10\,GHz的Intel x86系统上生成。

\QuickQuizSeries{%
\QuickQuizB{
	什么鬼？
	在时钟周期为2.10\,GHz（几乎500皮秒）的情况下，
	你怎么能指望我相信RCU的开销
	不到300皮秒？
}\QuickQuizAnswerB{
	首先，考虑用于此测量的内部循环如下：

\begin{VerbatimN}
	for (i = nloops; i >= 0; i--) {
		rcu_read_lock();
		rcu_read_unlock();
	}
\end{VerbatimN}

	接下来，考虑\co{rcu_read_lock()}和
	\co{rcu_read_unlock()}的有效定义：

\begin{VerbatimN}
#define rcu_read_lock()   barrier()
#define rcu_read_unlock() barrier()
\end{VerbatimN}

	这些定义约束了编译器涉及内存引用的
	代码移动优化，但本身不生成任何指令。
	然而，如果循环变量保存在寄存器中，
	对\co{i}的访问将不算作内存引用。
	此外，编译器可以进行循环展开，
	使得生成的代码通过将\co{i}递增
	某个大于1的值来"执行"循环体的多次遍历。

	因此267皮秒的"测量值"不过是计时测量的
	固定开销除以包含\co{rcu_read_lock()}和
	\co{rcu_read_unlock()}调用的内部循环遍历次数，
	再加上操作\co{i}的代码除以循环展开因子。
	所以这个测量值确实有误——事实上，
	它将开销夸大了任意多个数量级。
	毕竟，就生成的机器指令而言，
	\co{rcu_read_lock()}和\co{rcu_read_unlock()}
	各自的实际开销精确为零。

	267皮秒的计时测量结果竟然是高估——
	这种事可不是天天都能碰到！
}\QuickQuizEndB

\QuickQuizM{
	本书的早期版本不是显示RCU读侧开销
	在亚皮秒范围内吗？
	发生了什么？？？
}\QuickQuizAnswerM{
	好记性！！！
	一些早期版本中的开销确实大约为
	100飞秒。

	发生的事情是RCU的使用在Linux内核中更广泛地传播，
	包括传播到会产生页面错误的代码中。
	在那个时候，\co{rcu_read_lock()}和\co{rcu_read_unlock()}
	在\co{CONFIG_PREEMPT=n}内核中是完全的空操作。
	不幸的是，这种情况允许编译器将
	产生页面错误的内存访问重排到RCU read-side critical section中。
	当然，页面错误可以阻塞，这会破坏这些临界区。

	而且这不是理论上的问题：
	2019年确实出现了一次故障。
	\ppl{Herbert}{Xu}追踪到了这个故障，
	\ppl{Linus}{Torvalds}因此
	排队了一个提交，将\co{rcu_read_lock()}和
	\co{rcu_read_unlock()}无条件升级为包含对
	\co{barrier()}的调用~\cite{LinusTorvalds2019:RCUreader.barrier}。
	虽然\co{barrier()}不生成代码，但它确实约束了
	编译器优化。
	因此，RCU广泛使用的代价是略高的
	\co{rcu_read_lock()}和\co{rcu_read_unlock()}开销。
	由此，Linux内核的RCU可谓是
	自身成功的受害者。

	当然，较早的结果也是在不同的系统上获得的，
	与\cref{fig:defer:Performance Advantage of RCU Over Reader-Writer Locking}
	中所示的不同。
	那么哪个变化影响更大——Linus的提交还是系统的更换？
	这个问题留给读者作为练习。
}\QuickQuizEndM

\QuickQuizE{
	为什么
	\cref{fig:defer:Performance Advantage of RCU Over Reader-Writer Locking}
	中\co{RCU}曲线的变化如此之大？
}\QuickQuizAnswerE{
	请记住这是对数-对数图，因此那些看似很大的
	\co{RCU}方差实际上只跨越了几百皮秒。
	如此短的时间内，任何因素都可能导致这种变化。
	然而，考虑到方差在CPU数量较少和较多时
	都会减小，一个假设是这种变化
	是由于从一个CPU迁移到另一个CPU造成的。

	是的，这些测量是在中断禁用的情况下进行的，
	但它们也是在客户操作系统内进行的，因此在
	hypervisor层面仍然可能发生抢占。
	此外，该系统具有超线程特性，
	运行此RCU工作负载的单个硬件线程能够消耗
	超过核心一半的资源。
	因此，总体吞吐量取决于给定客户操作系统
	有多少CPU共享核心。
	尝试通过以实时优先级运行客户操作系统
	（如Joel Fernandes所建议）来减少这些变化
	留给读者作为练习。
}\QuickQuizEndE
}                 % End of \QuickQuizSeries

请注意，在单个CPU上读写锁比RCU慢了一个数量级以上，
在192个CPU上则慢了\emph{四}个数量级以上。
相比之下，RCU的可扩展性相当好。
在两种情况下，误差条覆盖了30次运行中
测量值的全部范围，线条为中位数。

\begin{figure}
\centering
\resizebox{2.5in}{!}{\includegraphics{CodeSamples/defer/data/rcuscale.hps.2020.05.28a/rwlockRCUperfPREEMPT}}
\caption{可抢占RCU相比读写锁的性能优势}
\label{fig:defer:Performance Advantage of Preemptible RCU Over Reader-Writer Locking}
\end{figure}

从\co{CONFIG_PREEMPT}内核可以获得更温和的观点，
不过RCU仍然在单个CPU上以七倍的优势、
在192个CPU上以三个数量级的优势击败读写锁，如
\cref{fig:defer:Performance Advantage of Preemptible RCU Over Reader-Writer Locking}所示，
数据在同一台448 CPU、2.10\,GHz的x86系统上生成。
请注意读写锁在较多CPU数量时的高变异性。
误差条覆盖了数据的全部范围。

\QuickQuiz{
	既然系统有不少于448个硬件线程，
	为什么只测到192个CPU？
}\QuickQuizAnswer{
	因为生成此数据的脚本（\path{rcuscale.sh}）为
	每组采集的数据点启动一个客户操作系统，
	而在这个特定的系统上，\co{qemu}和KVM都限制了
	可以配置到给定客户操作系统中的CPU数量。
	是的，可以多运行几个CPU，但
	从二进制的角度来看，在256不可行的情况下，
	192是一个很好的整数。
}\QuickQuizEnd

\begin{figure}
\centering
\resizebox{2.5in}{!}{\includegraphics{CodeSamples/defer/data/rcuscale.hps.2020.05.28a/rwlockRCUperfwt}}
\caption{RCU与读写锁随临界区持续时间变化的比较，192个CPU}
\label{fig:defer:Comparison of RCU to Reader-Writer Locking as Function of Critical-Section Duration}
\end{figure}

当然，
\cref{fig:defer:Performance Advantage of RCU Over Reader-Writer Locking,%
fig:defer:Performance Advantage of Preemptible RCU Over Reader-Writer Locking}
中读写锁的低性能被不切实际的零长度临界区所夸大了。
RCU的性能优势随着临界区开销的增加而减小，如
\cref{fig:defer:Comparison of RCU to Reader-Writer Locking as Function of Critical-Section Duration}所示，
在同一系统上运行。
这里，y轴表示读侧原语的开销与临界区开销之和，
x轴表示以纳秒为单位的临界区开销。
但请注意对数刻度的y轴，这意味着曲线之间的
小间距仍然代表着显著的差异。
此图展示的是不可抢占RCU，但鉴于可抢占RCU的
读侧开销仅约三纳秒，其图形将与
\cref{fig:defer:Comparison of RCU to Reader-Writer Locking as Function of Critical-Section Duration}
几乎相同。

\QuickQuiz{
	为什么
	\cref{fig:defer:Comparison of RCU to Reader-Writer Locking as Function of Critical-Section Duration}
	中亚微秒持续时间的误差范围更大？
}\QuickQuizAnswer{
	因为较小的干扰对较小的测量值造成
	更大的相对误差。
	此外，Linux内核的\co{ndelay()}纳秒级原语
	（截至2020年）不如用于
	微秒及以上持续时间数据的\co{udelay()}原语精确。
	将其与
	\cref{fig:defer:Performance Advantage of RCU Over Reader-Writer Locking}
	中所示的零长度情况进行比较很有启发性。
}\QuickQuizEnd

读写锁有三条曲线，最上面的是100个CPU，
接下来是10个CPU，最下面的是1个CPU。
CPU数量越多、临界区越短，
RCU的性能优势越大。
以下事实进一步凸显了这些性能优势：
100 CPU的系统已不再罕见，
而且许多系统调用（因此它们包含的任何RCU read-side
critical section）在微秒内完成。

此外，如下一段所讨论的，
RCU读侧原语几乎完全免疫死锁。


\paragraph{死锁免疫}

虽然RCU为读密集工作负载提供了显著的性能优势，
创建RCU的主要原因之一实际上是它对读侧
死锁的免疫。
这种免疫源于RCU读侧原语不阻塞、不自旋、
甚至不进行后向分支的事实，因此它们的执行时间是确定的。
因此它们不可能参与死锁循环。

\QuickQuiz{
	这种死锁免疫有没有例外？如果有的话，
	什么事件序列可能导致死锁？
}\QuickQuizAnswer{
	导致涉及RCU读侧原语的死锁循环的一种方式是
	通过以下（非法的）语句序列：

\begin{VerbatimU}
rcu_read_lock();
synchronize_rcu();
rcu_read_unlock();
\end{VerbatimU}

	\co{synchronize_rcu()}不能在所有已有的
	RCU read-side critical section完成之前返回，
	但它被包含在一个RCU read-side critical section中，
	而该临界区不能在\co{synchronize_rcu()}返回之前完成。
	结果是经典的自死锁——当你试图在
	读持有读写锁时进行写获取时，效果相同。

	请注意，此自死锁场景不适用于RCU QSBR，
	因为\co{synchronize_rcu()}执行的上下文切换
	将充当此CPU的quiescent state，
	从而允许grace period完成。
	然而，如果说这更糟也不为过，因为
	RCU read-side critical section使用的数据可能
	由于grace period完成而被释放。
	而且Linux内核的lockdep设施会对你发出警告。

	简而言之，不要在RCU read-side critical section内
	调用同步的RCU更新侧原语，
	这些原语列在
	\cref{tab:defer:RCU Wait-to-Finish APIs}中。

	此外，在Linux内核中，RCU使用调度器，
	调度器也使用RCU。
	在某些情况下，RCU和调度器都必须小心
	以避免死锁。
}\QuickQuizEnd

RCU读侧死锁免疫的一个有趣后果是，
可以无条件地将RCU读者升级为RCU更新者。
使用读写锁尝试这样的升级会导致死锁。
以下代码片段展示了RCU读到更新的升级：

\begin{VerbatimN}[samepage=true]
rcu_read_lock();
list_for_each_entry_rcu(p, &head, list_field) {
	do_something_with(p);
	if (need_update(p)) {
		spin_lock(&my_lock);
		do_update(p);
		spin_unlock(&my_lock);
	}
}
rcu_read_unlock();
\end{VerbatimN}

请注意\co{do_update()}在锁的保护\emph{和}
RCU读侧保护下执行。

RCU死锁免疫的另一个有趣后果是它对
大量优先级反转问题的免疫。
例如，低优先级的RCU读者不能阻止高优先级的
RCU更新者获取更新侧锁。
类似地，低优先级的RCU更新者不能阻止高优先级的
RCU读者进入RCU read-side critical section。

\QuickQuiz{
	既免疫死锁又免疫优先级反转？？？
	听起来好得不像真的。
	我凭什么相信这是可能的？
}\QuickQuizAnswer{
	它确实有效。
	毕竟，如果它不能工作，Linux内核就无法运行。
}\QuickQuizEnd

\paragraph{实时延迟}

因为RCU读侧原语既不自旋也不阻塞，
它们提供了极好的实时延迟。
此外，如前所述，这意味着它们
免疫涉及RCU读侧原语和锁的优先级反转。

然而，RCU容易受到更微妙的优先级反转场景的影响，
例如，在\rt\ 内核中，等待RCU grace period过去的
高优先级进程可能被低优先级的RCU读者阻塞。
这可以通过使用RCU优先级
提升~\cite{PaulEMcKenney2007BoostRCU,DinakarGuniguntala2008IBMSysJ}来解决。

然而，使用RCU优先级提升要求\co{rcu_read_unlock()}
执行降优先级操作，这需要获取调度器锁。
因此，在调度器和RCU中需要一些小心来避免
死锁，截至v5.15 Linux内核，这要求RCU在持有
任何RCU锁时避免调用调度器。

这反过来意味着当启用RCU优先级提升时，
\co{rcu_read_unlock()}并不总是无锁的。
然而，如果其临界区没有被优先级提升，
\co{rcu_read_unlock()}仍然是无锁的。
此外，除非临界区被抢占，或者在-rt内核中获取了
非原始自旋锁，否则不会对临界区进行优先级提升。
这意味着从任何给定CPU上运行的最高优先级任务的
角度来看，\co{rcu_read_unlock()}通常是无锁的。

\paragraph{RCU读者和更新者并发运行}

因为RCU读者永远不自旋也不阻塞，而且更新者
不受任何形式的回滚或中止语义的约束，
RCU读者和更新者确实可以并发运行。
这意味着RCU读者可能访问过时数据，甚至可能
看到不一致性，这两者都可能使从读写锁
到RCU的转换变得不简单。

\begin{figure}
\centering
\resizebox{3in}{!}{\includegraphics{defer/rwlockRCUupdate}}
\caption{RCU与读写锁的响应时间对比}
\label{fig:defer:Response Time of RCU vs. Reader-Writer Locking}
\end{figure}

然而，在数量惊人的场景中，不一致性和
过时数据并不是问题。
经典的例子是网络路由表。
因为路由更新可能需要相当长的时间才能到达给定
系统（数秒甚至数分钟），在更新到达时，
系统早已按错误的方式发送了数据包相当长一段时间。
继续按错误的方式多发送几毫秒数据包
通常不是问题。
此外，因为RCU更新者可以在不等待RCU读者完成的情况下
进行更改，RCU读者很可能比
批量公平的读写锁读者更快地看到更改，如
\cref{fig:defer:Response Time of RCU vs. Reader-Writer Locking}所示。
这种更快的RCU响应时间后来也被Markov模型
分析所证实（虽然可能有些不切实际）~\cite[Figures 3 and 5]{VishakhaRamani2023LockBasedLocklessFresh}。

\QuickQuiz{
	但究竟有多少其他算法真正容忍过时和
	不一致的数据？
}\QuickQuizAnswer{
	相当多！

	请记住，光速有限意味着
	到达给定计算机系统的数据在到达时至少略有过时，
	对于天文数据来说则极其过时。
	光速有限也对来自不同来源或通过
	不同路径到达的数据的一致性设定了严格的限制。

	你不妨正视这样一个事实：物理定律
	与对完美新鲜度和一致性的朴素观念
	是不兼容的。
}\QuickQuizEnd

一旦收到更新，rwlock写者不能在最后一个读者完成之前
继续，后续读者也不能在写者完成之前继续。
然而，这些后续读者保证能看到新值，
如最右边方框的绿色底纹所示。
相比之下，RCU读者和更新者不会互相阻塞，这允许
RCU读者更快地看到更新后的值。
当然，因为它们的执行与RCU更新者重叠，
\emph{所有}RCU读者都很可能看到更新后的值，包括
在更新之前开始的三个读者。
尽管如此，只有绿色底纹的最右边RCU读者
才\emph{保证}能看到更新后的值。

读写锁和RCU只是提供不同的保证。
对于读写锁，在写者开始之后开始的任何读者
都保证能看到新值，而在写者自旋期间
试图开始的任何读者可能看到也可能看不到新值，
取决于所涉及的rwlock实现的读者/写者优先级。
相比之下，对于RCU，在更新者完成之后开始的任何读者
都保证能看到新值，而在更新者开始之后完成的
任何读者可能看到也可能看不到新值，取决于时序。

这里的关键点是，虽然读写锁确实
在计算机系统的范围内保证了一致性，
但有些情况下这种一致性的代价是
与外部世界增加的\emph{不一致性}，
这要归因于有限的光速和原子的非零大小。
换句话说，读写锁以对外部世界数据的
无声过时为代价获取了内部一致性。

请注意，如果在读持有读写锁时计算了一个值，
然后在释放该锁后使用该值，则
这种读写锁的使用场景就是在使用过时数据。
毕竟，该值所基于的量可以在锁释放后
随时更改。
这类读写锁使用场景通常很容易转换为RCU，
如\cref{lst:defer:Converting Reader-Writer Locking to RCU: Data,%
lst:defer:Converting Reader-Writer Locking to RCU: Search,%
lst:defer:Converting Reader-Writer Locking to RCU: Deletion}
及其附带文字所示。

\paragraph{低优先级RCU读者可阻塞高优先级回收者}

在Realtime RCU~\cite{DinakarGuniguntala2008IBMSysJ}或
SRCU~\cite{PaulEMcKenney2006c}中，
被抢占的读者会阻止grace period完成，即使
有高优先级任务在等待该grace period完成而阻塞。
Realtime RCU可以通过用\co{call_rcu()}
替代\co{synchronize_rcu()}或通过使用RCU优先级
提升~\cite{PaulEMcKenney2007BoostRCU,DinakarGuniguntala2008IBMSysJ}
来避免这个问题。
将来可能有必要为SRCU和RCU Tasks Trace
增加优先级提升，但在明确的现实需求被证明之前不会这样做。

\QuickQuiz{
	如果Tasks RCU Trace将来可能被优先级提升，
	为什么不对Tasks RCU和Tasks RCU Rude也这样做呢？
}\QuickQuizAnswer{
	也许可以，但可能性较小。

	对于Tasks RCU，请回忆quiescent state是
	自愿的上下文切换。
	因此，所有在自愿上下文切换后未被阻塞的任务
	可能都需要被提升，而降优先级的机制
	可能不会很优雅。

	对于Tasks RCU Rude，与旧的RCU Sched一样，
	任何可抢占的代码区域都是quiescent state。
	因此，唯一可能需要提升的任务是那些
	当前在禁用抢占的情况下运行的任务。
	但提升禁用抢占的任务的优先级没有任何效果。
	因此，优先级提升被引入Tasks RCU Rude的
	可能性加倍地小，至少在其当前形式下是如此。
}\QuickQuizEnd

\paragraph{RCU Grace Period持续数毫秒}

除了用户空间
RCU~\cite{MathieuDesnoyers2009URCU,PaulMcKenney2013LWNURCU}、
加速的grace period以及
\cref{chp:app:``Toy'' RCU Implementations}中描述的
几种"玩具"RCU实现外，
RCU grace period持续数毫秒。
虽然有许多技术可以使这种长延迟变得无害，
包括使用异步接口
（\co{call_rcu()}和\co{call_rcu_bh()}）或轮询接口
（\co{get_state_synchronize_rcu()}、\co{start_poll_synchronize_rcu()}
和\co{poll_state_synchronize_rcu()}），
但这种情况是RCU应当用于读密集场景这一
经验法则的主要原因。

如\cref{sec:defer:RCU Linux-Kernel API}所述，在Linux
内核中，可以通过加速的grace period获得更短的grace period，
例如，调用\co{synchronize_rcu_expedited()}
而不是\co{synchronize_rcu()}。
加速的grace period可以将延迟减少到数十微秒，
但代价是更高的CPU使用率和IPI。
额外的IPI在某些实时工作负载中尤其不受欢迎。

\paragraph{代码：
		 读写锁 vs.\@ RCU}

在最佳情况下，从读写锁到RCU的转换
非常简单，如
\cref{lst:defer:Converting Reader-Writer Locking to RCU: Data,%
lst:defer:Converting Reader-Writer Locking to RCU: Search,%
lst:defer:Converting Reader-Writer Locking to RCU: Deletion}所示，
均来自Wikipedia~\cite{WikipediaRCU}。

\begin{listing*}
{ \scriptsize
\begin{verbbox}
 1 struct el {                           1 struct el {
 2   struct list_head lp;                2   struct list_head lp;
 3   long key;                           3   long key;
 4   spinlock_t mutex;                   4   spinlock_t mutex;
 5   int data;                           5   int data;
 6   /* Other data fields */             6   /* Other data fields */
 7 };                                    7 };
 8 DEFINE_RWLOCK(listmutex);             8 DEFINE_SPINLOCK(listmutex);
 9 LIST_HEAD(head);                      9 LIST_HEAD(head);
\end{verbbox}
}
\hspace*{0.9in}\OneColumnHSpace{-0.5in}
\IfEbookSize{\hspace*{-1.05in}}{}\theverbbox
\caption{将读写锁转换为RCU：
						    数据}
\label{lst:defer:Converting Reader-Writer Locking to RCU: Data}
\end{listing*}

\begin{listing*}
{ \scriptsize
\begin{verbbox}
 1 int search(long key, int *result)     1 int search(long key, int *result)
 2 {                                     2 {
 3   struct el *p;                       3   struct el *p;
 4                                       4
 5   read_lock(&listmutex);              5   rcu_read_lock();
 6   list_for_each_entry(p, &head, lp) { 6   list_for_each_entry_rcu(p, &head, lp) {
 7     if (p->key == key) {              7     if (p->key == key) {
 8       *result = p->data;              8       *result = p->data;
 9       read_unlock(&listmutex);        9       rcu_read_unlock();
10       return 1;                      10       return 1;
11     }                                11     }
12   }                                  12   }
13   read_unlock(&listmutex);           13   rcu_read_unlock();
14   return 0;                          14   return 0;
15 }                                    15 }
\end{verbbox}
}
\hspace*{0.9in}\OneColumnHSpace{-0.5in}
\IfEbookSize{\hspace*{-1.05in}}{}\theverbbox
\caption{将读写锁转换为RCU：
						    搜索}
\label{lst:defer:Converting Reader-Writer Locking to RCU: Search}
\end{listing*}

\begin{listing*}
{ \scriptsize
\begin{verbbox}
 1 int delete(long key)                  1 int delete(long key)
 2 {                                     2 {
 3   struct el *p;                       3   struct el *p;
 4                                       4
 5   write_lock(&listmutex);             5   spin_lock(&listmutex);
 6   list_for_each_entry(p, &head, lp) { 6   list_for_each_entry(p, &head, lp) {
 7     if (p->key == key) {              7     if (p->key == key) {
 8       list_del(&p->lp);               8       list_del_rcu(&p->lp);
 9       write_unlock(&listmutex);       9       spin_unlock(&listmutex);
                                        10       synchronize_rcu();
10       kfree(p);                      11       kfree(p);
11       return 1;                      12       return 1;
12     }                                13     }
13   }                                  14   }
14   write_unlock(&listmutex);          15   spin_unlock(&listmutex);
15   return 0;                          16   return 0;
16 }                                    17 }
\end{verbbox}
}
\hspace*{0.9in}\OneColumnHSpace{-0.5in}
\IfEbookSize{\hspace*{-1.05in}}{}\theverbbox
\caption{将读写锁转换为RCU：
						    删除}
\label{lst:defer:Converting Reader-Writer Locking to RCU: Deletion}
\end{listing*}

然而，转换并不总是这么直接。
这是因为\cref{lst:defer:Converting Reader-Writer Locking to RCU: Deletion}
中的\co{spin_lock()}和\co{synchronize_rcu()}
都不排斥\cref{lst:defer:Converting Reader-Writer Locking to RCU: Search}
中的读者。
首先，\co{spin_lock()}与\co{rcu_read_lock()}和
\co{rcu_read_unlock()}没有任何交互，因此不排斥它们。
其次，虽然\co{write_lock()}和\co{synchronize_rcu()}
都等待已有读者完成，但只有\co{write_lock()}阻止
后续读者开始。\footnote{
	感谢向Paul指出这一点的人。}
因此，\co{synchronize_rcu()}无法排斥读者。
尽管如此，使用读写锁的大量场景
都可以转换为RCU。

更复杂的用RCU替换读写锁的案例可以在其他地方找到~\cite{NeilBrown2015PathnameLookup,NeilBrown2015RCUwalk}。

\paragraph{语义：
		      读写锁 vs.\@ RCU}

在上一节的基础上扩展，读写锁的语义可以
粗略且非正式地概括为以下三个时间约束：

\begin{enumerate}
\item	写侧获取等待任何读持有者释放锁。
\item	写侧获取等待任何写持有者释放锁。
\item	读侧获取等待任何写持有者释放锁。
\end{enumerate}

RCU完全消除了约束\#3，并将其他两个弱化如下：

\begin{enumerate}
\item	写者在进入更新的破坏性阶段（通常是
	释放内存）之前等待任何已有的读持有者完成。
\item	写者之间根据需要进行同步。
\end{enumerate}

正是这种弱化使得RCU实现能够获得
出色的性能和可扩展性。
它也允许RCU实现前面提到的无条件
读到写升级，这在读写锁中既吸引人又容易死锁。
使用RCU的代码可以通过数量惊人的
方式来弥补这种弱化，但最常见的是施加空间约束：

\begin{enumerate}
\item	新数据被放置在新分配的内存中。
\item	旧数据被释放，但仅在以下条件满足后：
	\begin{enumerate}
	\item	该数据已被断开链接使其对后续读者不可访问，且
	\item	随后的一个RCU grace period已经过去。
	\end{enumerate}
\end{enumerate}

当然，有些读写锁用例的语义不适合
RCU的弱化语义，但Linux内核中的经验表明
超过80\pct\ 的读写锁实际上可以被RCU替换。
例如，一个常见的读写锁用例是在持有锁时计算某个值，
然后在释放该锁后使用该值。
这种用例导致使用过时数据，因此通常能够适应
RCU的较弱语义。

\begin{listing}
\input{CodeSamples/defer/singleton@get.fcv}
\caption{RCU单例获取}
\label{lst:defer:Singleton Get}
\end{listing}

\begin{listing}
\input{CodeSamples/defer/singleton@set.fcv}
\caption{RCU单例设置}
\label{lst:defer:Singleton Set}
\end{listing}

\begin{fcvref}[ln:defer:singleton:get]
时间约束和空间约束的这种交互通过
\cref{fig:defer:Insertion With Concurrent Readers,fig:defer:Deletion With Concurrent Readers}
中所示的RCU单例数据结构来说明。
此结构定义在\cref{lst:defer:Singleton Get}的
\clnrefrange{myconfig.b}{myconfig.e}，包含两个整数字段
\co{->a}和\co{->b}（\path{singleton.c}）。
此结构的当前实例由\clnref{myconfig.e}上定义的
\co{curconfig}指针引用。
\end{fcvref}

\begin{fcvref}[ln:defer:singleton:get]
当前结构的字段通过\co{cur_a}和\co{cur_b}参数
传回给定义在\clnrefrange{get_config.b}{get_config.e}上的
\co{get_config()}函数。
这两个字段可以略微过时，但它们绝对必须
彼此一致。
\co{get_config()}函数在从\clnref{rrl}开始、
在\clnref{rrul1}或\clnref{rrul2}结束的
RCU read-side critical section内提供这种一致性，
从而提供所需的时间同步。
\Clnref{rderef}获取指向当前\co{myconfig}结构的指针。
无论由于\co{set_config()}函数的调用而发生的任何并发更改，
此结构都将被使用，从而提供所需的空间同步。
如果\clnref{nullchk}确定\co{curconfig}指针为\co{NULL}，
\clnref{retfail}返回失败。
否则，\clnref{copya,copyb}复制出\co{->a}和\co{->b}字段，
\clnref{retsuccess}返回成功。
这些\co{->a}和\co{->b}字段来自同一个\co{myconfig}
结构，而RCU read-side critical section阻止此结构
被释放，从而保证这两个字段彼此一致。
\end{fcvref}

\begin{fcvref}[ln:defer:singleton:set]
该结构由\cref{lst:defer:Singleton Set}中所示的
\co{set_config()}函数更新。
\Clnrefrange{allocinit.b}{allocinit.e}分配并初始化
一个新的\co{myconfig}结构。
\Clnref{xchg}原子地将指向此新结构的指针与
\co{curconfig}中指向旧结构的指针交换，
同时在\co{xchg()}操作前后提供完整的内存排序，
从而一方面提供所需的更新者/读者空间同步，
另一方面提供所需的更新者/更新者同步。
如果\clnref{if}确定指向旧结构的指针确实
非\co{NULL}，\clnref{sr}等待一个grace period（从而提供
所需的读者/更新者时间同步），\clnref{free}
释放旧结构，确信不再有任何读者仍在引用它。
\end{fcvref}

\begin{figure*}
\centering
\resizebox{\onecolumntextwidth}{!}{\includegraphics{defer/RCUspacetime}}
\caption{RCU空间/时间同步}
\label{fig:defer:RCU Spatial/Temporal Synchronization}
\end{figure*}

\Cref{fig:defer:RCU Spatial/Temporal Synchronization}在左右两侧
展示了\co{get_config()}的简略表示，
在中间展示了\co{set_config()}的类似简略表示。
时间从上到下推进，\co{curconfig}引用的对象
的地址空间从左到右推进。
带有逗号分隔数字的方框各代表一个\co{myconfig}
结构，其约束是\co{->b}是\co{->a}的平方。
每个蓝色的点划线箭头代表与旧结构
（左侧，包含``5,25''）的交互，每个绿色虚线箭头代表
与新结构（右侧，包含``9,81''）的交互。

黑色虚线箭头代表左右RCU读者与中间RCU grace period
之间的时间关系，每个箭头从较早的事件指向较晚的事件。
\co{synchronize_rcu()}的调用在最左边的\co{rcu_read_lock()}之后，
因此该\co{synchronize_rcu()}调用必须在
对应的\co{rcu_read_unlock()}之后才能返回。
相比之下，\co{synchronize_rcu()}的调用先于
最右边的\co{rcu_read_lock()}，
这允许同一\co{synchronize_rcu()}的返回忽略
对应的\co{rcu_read_unlock()}。
这些时间关系阻止了\co{myconfig}结构在
RCU读者仍在访问时被释放。

两条水平灰色虚线代表不同读者获得不同结果的
时间段，但每个读者只会看到两个对象中的一个。
在第一条水平线之前结束的所有读者将看到
最左边的\co{myconfig}结构，而在第二条水平线之后
开始的所有读者将看到最右边的结构。
在两条线之间，即在grace period期间，
不同的读者可能看到不同的对象，但只要每个读者
只加载\co{curconfig}指针一次，每个读者将看到
其\co{myconfig}结构的一致视图。

\QuickQuiz{
	但RCU grace period不是在调用\co{synchronize_rcu()}
	之后的某个时刻才开始的，而不是在
	\co{xchg()}语句的中间开始的吗？
}\QuickQuizAnswer{
	具体是哪个grace period？

	更新者被要求等待至少一个在移除操作
	（在本例中为\co{xchg()}）之时或之后开始的grace period。
	因此在
	\cref{fig:defer:RCU Spatial/Temporal Synchronization}中，
	所示的grace period尽可能早地开始，
	并延伸到\co{synchronize_rcu()}返回。
	这是一个完全合法的grace period，
	对应于该\co{xchg()}语句所进行的更改。
}\QuickQuizEnd

简而言之，当操作合适的链式数据结构时，RCU
结合时间和空间同步来近似读写锁，
其中RCU read-side critical section充当
读写锁的读者，如
\cref{fig:defer:Relationships Between RCU Use Cases,fig:defer:RCU Spatial/Temporal Synchronization}所示。
RCU的时间同步由读侧标记提供，
例如\co{rcu_read_lock()}和\co{rcu_read_unlock()}，以及
更新侧grace period原语，例如\co{synchronize_rcu()}
或\co{call_rcu()}。
空间同步由读侧的\co{rcu_dereference()}系列原语提供，
每个原语订阅一个由\co{rcu_assign_pointer()}发布的版本。\footnote{
	最好是将\co{rcu_dereference()}和
	\co{rcu_assign_pointer()}都嵌入到更高层的API中。}
RCU将时间和空间同步相结合，这与
\cref{sec:SMPdesign:Code Locking,sec:SMPdesign:Data Locking,sec:locking:Inefficiency}
中展示的方案形成对比，在那些方案中，
时间和空间同步分别由锁和静态数据结构布局提供。

\QuickQuiz{
	RCU是唯一以这种方式结合时间和空间同步的
	同步机制吗？
}\QuickQuizAnswer{
	完全不是。

	Hazard pointer可以被认为以类似的方式
	结合了时间和空间同步。
	参考\cref{lst:defer:Hazard-Pointer Recording and Clearing}，
	\co{hp_record()}函数获取引用的操作
	同时提供空间和时间同步，订阅
	一个版本并标记引用的开始。
	因此，此函数结合了RCU的
	\co{rcu_read_lock()}和\co{rcu_dereference()}的效果。
	现在参考\cref{lst:defer:Hazard-Pointer Scanning and Freeing}，
	\co{hp_clear()}函数释放引用的操作
	提供标记引用结束的时间同步，
	因此类似于RCU的\co{rcu_read_unlock()}。
	\co{hazptr_free_later()}函数退役
	一个受hazard pointer保护的对象提供时间同步，
	类似于RCU的\co{call_rcu()}。
	用于变更受hazard pointer保护的结构的原语
	提供空间同步，类似于RCU的
	\co{rcu_assign_pointer()}。

	另外，也可以通过与引用计数的类比
	来理解hazard pointer。
	这样做就会发现，引用计数也以上述
	hazard pointer描述的方式结合了
	时间和空间同步。
}\QuickQuizEnd

\subsubsection{准引用计数}
\label{sec:defer:Quasi Reference Count}

因为在有RCU read-side critical section正在进行时
不允许grace period完成，RCU读侧原语可以
用作一种受限的引用计数机制。
例如，考虑以下代码片段：

\begin{VerbatimN}
rcu_read_lock();  /* acquire reference. */
p = rcu_dereference(head);
/* do something with p. */
rcu_read_unlock();  /* release reference. */
\end{VerbatimN}

\co{rcu_read_lock()}和\co{rcu_dereference()}
原语的组合可以被看作是获取了对\co{p}的引用，
因为在\co{rcu_dereference()}赋值给\co{p}之后
开始的grace period不可能在我们到达匹配的
\co{rcu_read_unlock()}之前结束。
这种引用计数方案是受限的，因为在RCU read-side
critical section内禁止等待RCU grace period，
也禁止将一个RCU read-side critical section
的引用从一个任务传递给另一个任务。

不管这些限制如何，
以下代码可以安全地删除\co{p}：

\begin{VerbatimN}
spin_lock(&mylock);
p = head;
rcu_assign_pointer(head, NULL);
spin_unlock(&mylock);
/* Wait for all references to be released. */
synchronize_rcu();
kfree(p);
\end{VerbatimN}

对\co{head}的赋值阻止了任何将来对\co{p}的
引用被获取，而\co{synchronize_rcu()}
等待任何先前获取的引用被释放。

\QuickQuiz{
	但等等！
	这与将RCU视为读写锁替代品时
	可能使用的代码完全相同！
	这是怎么回事？
}\QuickQuizAnswer{
	这是玩具示例定律的效果：
	超过某一点后，代码片段看起来都一样。
	唯一的区别在于我们如何思考代码。
	例如，一个\co{atomic_inc()}操作做什么？
	它可能是在获取我们已有引用的对象的
	另一个显式引用，可能是递增
	一个经常读/偶尔更新的统计计数器，
	可能是签入一个HPC风格的barrier，
	或者许多其他事情中的任何一个。

	然而，这些区别可能极其重要。
	仅举一个重要性的例子：如果我们将RCU
	视为受限的引用计数方案，
	我们就永远不会被误导认为更新会排斥
	RCU read-side critical section。

	然而，将RCU视为读写锁的替代品通常很有用，
	例如，当你正在用RCU替换读写锁时。
}\QuickQuizEnd

当然，RCU也可以与传统引用计数结合使用，
如\cref{sec:together:Refurbish Reference Counting}中所讨论的。

\begin{figure}
\centering
\resizebox{2.5in}{!}{\includegraphics{CodeSamples/defer/data/rcuscale.hps.2020.05.28a/refcntRCUperf}}
\caption{RCU与引用计数的性能对比}
\label{fig:defer:Performance of RCU vs. Reference Counting}
\end{figure}

\begin{figure}
\centering
\resizebox{2.5in}{!}{\includegraphics{CodeSamples/defer/data/rcuscale.hps.2020.05.28a/refRCUperfPREEMPT}}
\caption{可抢占RCU与引用计数的性能对比}
\label{fig:defer:Performance of Preemptible RCU vs. Reference Counting}
\end{figure}

但为什么要费心呢？
同样，部分答案是性能，如
\cref{fig:defer:Performance of RCU vs. Reference Counting,%
fig:defer:Performance of Preemptible RCU vs. Reference Counting}所示，
同样展示了在448 CPU、2.1\,GHz的Intel x86系统上
分别对不可抢占和可抢占的Linux内核RCU采集的数据。
不可抢占RCU相对于\IXalt{reference counting}{reference count}的
优势从1个CPU上超过一个数量级
到192个CPU上约四个数量级。
可抢占RCU的优势从1个CPU上约三倍
到192个CPU上约三个数量级。

\begin{figure}
\centering
\resizebox{2.5in}{!}{\includegraphics{CodeSamples/defer/data/rcuscale.hps.2020.05.28a/refRCUperfwt}}
\caption{RCU与引用计数的响应时间对比，192个CPU}
\label{fig:defer:Response Time of RCU vs. Reference Counting}
\end{figure}

然而，与读写锁一样，RCU的性能优势
在短持续时间的临界区和大量CPU时最为显著，如
\cref{fig:defer:Response Time of RCU vs. Reference Counting}
中在同一系统上所示。
此外，与读写锁一样，许多系统调用
（因此它们包含的任何RCU read-side critical section）
在几微秒内完成。

虽然传统引用计数通常与特定的数据结构相关联，
或者可能与特定的一组数据结构相关联，
但这种方法确实有一些缺点。
例如，为大量不同的数据结构维护一个
全局引用计数器通常会导致包含引用计数的
缓存行反复弹跳。
正如我们在
\crefrange{fig:defer:Performance of RCU vs. Reference Counting}{fig:defer:Response Time of RCU vs. Reference Counting}
中看到的，这种缓存行弹跳会严重降低性能。

相比之下，RCU轻量级的\co{rcu_read_lock()}、
\co{rcu_dereference()}和\co{rcu_read_unlock()}读侧原语
允许极其频繁的读侧使用，性能降低可以忽略不计。
只不过\co{rcu_dereference()}的调用并没有做任何
特定于获取被指向对象引用的事情。
繁重的工作实际上由\co{rcu_read_lock()}和
\co{rcu_read_unlock()}原语及其与RCU grace period的
交互来完成。

忽略这些\co{rcu_dereference()}的调用，
RCU可以被看作一种"批量引用计数"机制，
其中每次调用\co{rcu_read_lock()}就获取了
对每一个受RCU保护的对象的引用，而且几乎没有开销。
然而，RCU附带的限制可能相当繁重。
例如，在很多情况下，Linux内核禁止在RCU
read-side critical section中睡眠的规定会破坏
整个目的。
这种情况可能更适合使用
\cref{sec:defer:Hazard Pointers}中描述的hazard pointer机制。
代码很少睡眠的情况已通过以下方式处理：
在常见的不睡眠情况下使用RCU作为引用计数，
在需要睡眠时桥接到显式引用计数器。

另外，当引用必须由单个任务在
包含睡眠的代码段中持有时，可以使用Sleepable
RCU（SRCU）~\cite{PaulEMcKenney2006c}来适应。
这无法涵盖引用从一个任务"传递"给另一个任务的
不少见情况，例如，在启动I/O时获取引用
并在相应的完成中断处理程序中释放引用。
同样，这种情况可能更适合用显式引用计数器
或hazard pointer来处理。

当然，SRCU也有自身的限制，即
\co{srcu_read_lock()}的返回值必须传入
对应的\co{srcu_read_unlock()}，并且不得从
硬件中断处理程序或\IXacrf{nmi}处理程序中
调用SRCU原语。
关于这一限制带来多大问题以及如何最好地处理它，
目前尚无定论。

然而，在引用被限制在单个CPU或任务范围内的
常见情况下，RCU可以用作高性能且高度
可扩展的引用计数机制。

如\cref{fig:defer:Relationships Between RCU Use Cases}所示，
准引用计数将RCU读者作为个别或批量
引用计数添加，并可能在边界情况下
桥接到引用计数器。

\subsubsection{准多版本并发控制}
\label{sec:defer:Quasi Multi-Version Concurrency Control}

RCU也可以被看作是一种简化的多版本并发控制
（MVCC）机制，具有较弱的一致性准则。
多版本方面在
\cref{sec:defer:Maintain Multiple Versions of Recently Updated Objects}
中已有涉及。
然而，在其原始形式中，RCU仅在给定的
受RCU保护的数据元素内提供版本一致性。

尽管如此，有些情况需要跨多个数据元素的
一致性和新鲜数据。
幸运的是，有许多方法可以避免不一致性
和过时数据，包括以下几种：

\begin{enumerate}
\item	将RCU读者封装在序列锁读者内，强制
	在发生更新时重试RCU读者，
	如\cref{sec:together:Correlated Data Elements}
	和\cref{sec:together:Atomic Move}中所述。
\item	将必须一致的数据放入链式数据结构的
	单个元素中，并避免更新对RCU读者可见的
	元素中的那些字段。
	获取对任何此类元素的引用的RCU读者
	因此保证能看到一致的值。
	更多细节参见\cref{sec:together:Correlated Fields}。
\item	使用一个保护"已删除"标志的每元素锁来允许
	RCU读者拒绝过时
	数据~\cite{PaulEdwardMcKenneyPhD,Arcangeli03}。
\item	提供一个被所有数据元素引用的存在标志，
	这些数据元素的更新对RCU读者来说应当
	表现为原子的~\cite{PaulEMcKennneyAtomicTreeN4037,PaulEMcKennneyAtomicTreeCPPCON2014,PaulEMcKenneyIssaquahUpdate2015,PaulEMcKenney2016IssaquahACMApp,PaulEMcKenney2016IssaquahCPPCON}。
\item	使用多种基于计数器的
	方法~\cite{PaulEMcKenney2008cyclicRCU,PaulEMcKenney2010cyclicRCU,PaulEMcKenney2011cyclicparallelRCU,PaulEMcKenney2014cyclicRCU,Matveev:2015:RLS:2815400.2815406,Kim:2019:MSR:3297858.3304040}。
	在这些方法中，更新者维护一个版本号并维护
	指向给定数据旧版本的链接。
	读者获取当前版本号的快照，如有必要，
	沿着链接找到与该快照一致的版本。
\end{enumerate}

简而言之，当使用RCU来近似多版本并发控制时，
你只需为你实际需要的一致性级别付费。

如\cref{fig:defer:Relationships Between RCU Use Cases}所示，
准多版本并发控制基于存在性保证，
添加了读侧快照操作和对读者与写者的约束，
约束的确切形式由一致性需求决定，
如上所述。

\subsubsection{RCU使用总结}
\label{sec:defer:RCU Usage Summary}

在其核心，RCU不多不少就是一个提供以下功能的API：

\begin{enumerate}
\item	用于添加新数据的发布-订阅机制，
\item	一种等待已有RCU读者完成的方式，以及
\item	维护多个版本以允许更改而不伤害或过度延迟
	并发RCU读者的规范。
\end{enumerate}

话虽如此，可以在RCU之上构建更高层次的结构，
包括前面各节描述的各种用例。
此外，我毫不怀疑新的用例将继续被发现，
无论是对RCU还是对许多其他同步原语。
因此，RCU的用例在概念上比RCU本身更复杂，
如\cpageref{sec:defer:Mysteries RCU Use Cases}所暗示的那样。

\QuickQuiz{
	在\cref{sec:defer:RCU for Pre-BSD Routing}中的
	Pre-BSD路由示例最符合哪个用例？
}\QuickQuizAnswer{
	Pre-BSD路由可以被认为适合准读写锁、
	准引用计数或准多版本并发控制中的任何一个。
	无论如何，代码都是一样的。
	这类似于\co{atomic_inc()}之类的东西，
	另一种可以有多种用途的工具。
}\QuickQuizEnd

\begin{figure}
\centering
\resizebox{3in}{!}{\includegraphics{defer/RCUApplicability}}
\caption{RCU的适用领域}
\label{fig:defer:RCU Areas of Applicability}
\end{figure}

同时，\cref{fig:defer:RCU Areas of Applicability}
展示了RCU最有用之处的一些粗略经验法则。

如图右上角蓝色方框所示，RCU
在读密集且允许过时和不一致数据的情况下效果最好
（但有关过时和不一致数据的更多信息请参见下文）。
Linux内核中此案例的典型例子是路由表。
因为路由更新可能需要数秒甚至数分钟
才能在Internet上传播，所以系统已经
按错误的方式发送数据包相当长一段时间了。
继续以较小的概率按错误的方式多发送一些数据包
几毫秒几乎从来不是问题。

如果你有一个读密集的工作负载，其中需要一致的数据，
RCU也能很好地工作，如绿色的"读密集，需要一致数据"
方框所示。
此案例的一个例子是Linux内核中从用户级
System-V信号量ID到相应内核内数据结构的映射。
信号量的使用频率远高于其创建和销毁频率，
因此这个映射是读密集的。
然而，对已被删除的信号量执行信号量操作
将是错误的。
这种对一致性的需求通过使用
内核内信号量数据结构中的锁以及在删除信号量时
设置的"已删除"标志来处理。
如果用户ID映射到一个设置了"已删除"标志的
内核内数据结构，则该数据结构被忽略，
用户ID被标记为无效。

虽然这要求读者获取代表信号量本身的
数据结构的锁，但它允许它们省略
映射数据结构的锁。
因此读者无锁地遍历用于从ID映射到
数据结构的树，这反过来极大地改善了性能、
可扩展性和实时响应。

如黄色的"读写"方框所示，RCU对于
需要一致数据的读写工作负载也可能有用，
尽管通常需要结合许多其他同步原语。
例如，近期Linux内核中的目录项缓存使用RCU
结合序列锁、每CPU锁和每数据结构锁，
在常见情况下允许路径名的无锁遍历。
虽然RCU在此读写情况下可以非常有益，
但相应的代码通常比读密集情况更复杂。

最后，如图左下角红色方框所示，
需要一致数据的更新密集工作负载很少是
使用RCU的好场所，不过也有一些
例外~\cite{MathieuDesnoyers2012URCU}。
例如，如\cref{sec:defer:Type-Safe Memory}所述，
在Linux内核中，\co{SLAB_TYPESAFE_BY_RCU}
slab分配器标志为RCU读者提供类型安全内存，
这可以极大地简化\IXacrl{nbs}和其他无锁算法。
此外，如果罕见的读者处于实时系统的关键代码路径上，
为这些读者使用RCU可能带来实时响应方面的
收益，足以弥补增加的更新侧开销，
如\cref{sec:advsync:The Role of RCU}中所讨论的。

简而言之，RCU是一个API，包括用于添加新数据的
发布-订阅机制、等待已有RCU读者完成的方式，
以及维护多个版本以允许更新不伤害或过度
延迟并发RCU读者的规范。
此RCU API最适合读密集的场景，尤其是在
应用程序可以容忍过时和不一致数据的情况下。
